\input{../../stock/en-tete_v6.tex}

\newcommand{\titrechapitre}{Assembleur}
\newcommand{\numerochapitre}{2}

\renewcommand{\headrulewidth}{0.4pt}
\renewcommand{\footrulewidth}{0.4pt}
\fancyfoot[L]{Notes de Raphaël JONTEF}
\fancyfoot[C]{\thepage}
\fancyfoot[R]{Enseirb-Matmeca}
\fancyhead[L]{Cours de Mathématiques}
\fancyhead[R]{\titrechapitre}

\begin{document}
\input{../../stock/commands.tex}

% Page de titre custom
\setlength{\headheight}{15pt}
\begin{titlepage}
    \centering
    \vspace*{3cm}
    {\huge Chapitre \numerochapitre\par}
    {\Huge \bfseries\titrechapitre\par}
    \vspace{2cm}
    {\Large Notes rédigées par Raphaël JONTEF\par}
{\Large Avec l'aide de Github Copilot\par}
    \vspace{0.5cm}
    {\normalsize Cours enseigné par Amina Guermouche\par}
    \vspace{4cm}
    {\small Enseirb-Matmeca\par}
    {\small filière informatique\par}
    \vfill
    {\large \today\par}
\end{titlepage}

% plan du cours
\tableofcontents
\newpage


% -----------Corps du document--------------

\section{Introduction}

L'assembleur est le langage le plus bas niveau qui puisse exister, qui soit lisible par un être humain. 

\begin{definition}{}{Instruction Set Architecture}
    L'\notion{ensemble des instructions} (ou instruction set architecture, abrégé ISA) constitue l'ensemble des instructions que le programmeur peut utiliser.
\end{definition}

\subsection{Quand utiliser l'assembleur ?}

\begin{itemize}
    \item programmation de systèmes d'exploitation
\end{itemize}

\section{Instructions}

\subsection{Le \code{mov}}

\begin{definition}{}{instruction \code{mov}}
    \begin{lstASM}
        mov rax , 0
    \end{lstASM}
    Permet d'allouer 0 à la variable \code{x}.
    
\end{definition}

\begin{definition}{}{notation crochets}
    Les crochets \code{[ ]} permettent d'accéder à la valeur stockée à l'adresse mémoire contenue dans un registre. Par exemple~:
    \begin{lstASM}
        mov rax , [rbx]
    \end{lstASM}
    permet de copier dans \code{rax} la valeur stockée à l'adresse mémoire contenue dans \code{rbx}.

\end{definition}

\subsection{Le \code{add}}
\begin{definition}{}{instruction \code{add}}
    \begin{lstASM}
        add rax , rbx
    \end{lstASM}
    Permet d'ajouter la valeur de \code{rbx} à \code{rax} et de stocker le résultat dans \code{rax}.




\end{document}